% Reader question: how do Scott and Hayek motivate the chapter's design concern?
% Interpretive synthesis, not a theorem or a claim that all systems share one bottleneck.
% Evidence: sec:hayek; arrows link complementary cautions to a design proposal.
\begin{figure}[H]
\centering
\pdfigurehue{pdteal}%
\begin{tikzpicture}[pd figure,x=1cm,y=1cm]
  \node[pd state,minimum width=3.9cm,text width=3.5cm,align=center] (S)
    at (2.30,4.50) {Scott 1998: simplification can erase practical knowledge};
  \node[pd state,minimum width=3.9cm,text width=3.5cm,align=center] (H)
    at (7.10,4.50) {Hayek 1945: relevant knowledge can be dispersed and local};
  \node[pd focus state,minimum width=7.0cm,text width=6.4cm,align=center] (C)
    at (4.70,2.30) {shared design concern: a central summary may omit knowledge needed for action};
  \node[pd ready state,minimum width=7.0cm,text width=6.4cm,align=center] (P)
    at (4.70,0.30) {proposal: preserve evidence access and opportunities for practical judgment};
  \draw[pd focus arrow] (S) -- (C);
  \draw[pd focus arrow] (H) -- (C);
  \draw[pd ready arrow] (C) -- (P);
\end{tikzpicture}
\caption{Scott and Hayek offer complementary cautions: administrative legibility can erase local knowledge. The arrows motivate a design proposal; they do not prove a universal bottleneck or an effective intervention. [interpretive synthesis]}
\label{fig:hayek-scott}
\end{figure}
